% Nouveau EGA V (cote 25): archive cote 25, site page 34 (PDF page 35).
% typescript. Draft compared: none.
% AI-transcribed by Claude Opus 5.5 (claude-opus-5-5, adaptive thinking, effort high), Message Batches API, batch ega25b (msgbatch_01A7n5YuXubUPNK1zsepi8RG), 2026-11-13; not yet checked by eye.
\scanpage{25/34}
/irr $X_j$ de X telle que codim($T_i$,$X_j$)=k, et  = 0, (i.e. si f est quasi-fini \unsure{et k>0},  T n'ait pas de pt isolé x tel que $\dim_x X = k$ ).
\placedmargin{\unsure{30}}
Signalons le résultat dimensionnel suivant:

\uline{Proposition} 8.4. Soit T une partie fermée de X, et supposons
dodim(T,X) $\geq$ k. Alors pour tout \hand{$\xi\in\check{P}$}, on a dodim (\hand{$T_\xi$},$Y_{\hand{\xi}}$) $\geq$ k-1. Soit
Z l'ensemble des \hand{$\xi\in$} $\check{P}$ tels que codim ($T_{\hand{\xi}}$,$Y_{\hand{\xi}}$ ) = k-1, (i.e.
codim ($T_{\hand{\xi}}$,$Y_{\hand{\xi}}$) $<$ k) alors Z est une partie constructible rare de $\check{P}$ ,
i.e. $\overline{Z}$ est de codimension $\geq$ 1 dans $\check{P}$ . Pour qu'il soit de

codimension $\geq$ 2, il f et s que pour toute composante irréductible
$T_i$ de X de codimension égale à k,  et telle 
que dim f($T_i$) = 0, il existe une composante \hand{/}
La première assertion résulte aussitôt du lemme suivant (8.4.1.a), qui est

un remords au par.5 :

\uline{Lemme} 8.4.1. Soient X préschéma loc noeth, $\underline{L}$ un Module inversible sur
X , $\phi$ une section de $\underline{L}$, Y = V($\phi$), T une partie fermée de X. Supposons dodim(Y,X)$\geq$k. Alors :

a) Codim(T$\hand{\cap}$Y,Y) $\geq$ k-1.

b) Pour qu'on ait codim(T Y,Y) = k-1 i.e. codim(T$\hand{\cap}$Y,Y) $<$ k, il
f et s qu'il existe une composante irréductible $T_i$ de T, contenue dans
Y,  telle que codim($T_i$,X) = k, et telle que pour toute composante
irréductible $X_j$ de X  contenant $T_i$ et telle que dim $\underline{O}_{X_j,T_i}$ =
dim $\underline{O}_{X,T_i}$ ; (=k ), on ait $X_j \hand{\not\subset}$ Y.

La vérification de ce lemme est immédiate grâce aux généralités
de $O_{IV}$ sur la dimension. Sous les conditions de 8.4., on voit par 8.4.1. b) quels sont
les hyperplans exceptionnels $H_{\hand{\xi}}$ 
, Si
on excepte l'ensemble $Z_o$ des \hand{$\xi\in$} $\check{P}$ tels qu'il existe une composante
irr R de T ou de X telle dim f(R) $>$ 0 et que $R_{\hand{\xi}}$ soit de ``dimension trop grande'',
 (ensemble qui est de codimension $\geq$ 2 et par suite ne compte pas
les $H_{\hand{\xi}}$  exceptionnels sont ceux pour lesquels il existe
un $T_i$ avec codim($T_i$,X) = k , et  dim f(T ) = 0,

\mspage{34}
